\documentclass[12pt, a4paper]{article}

\usepackage[utf8]{inputenc}
\usepackage[T1]{fontenc}
\usepackage[ngerman]{babel}
\usepackage{amsmath, amssymb, amsthm}
\usepackage{geometry}
\usepackage{booktabs}
\usepackage{array}
\usepackage{hyperref}
\usepackage{xcolor}
\usepackage{titlesec}
\usepackage{enumitem}
\usepackage{lmodern}

\geometry{margin=2.5cm}

\hypersetup{
    colorlinks=true,
    linkcolor=blue!60!black,
    urlcolor=blue!60!black,
}

\titleformat{\section}{\large\bfseries}{}{0em}{}[\titlerule]
\titleformat{\subsection}{\normalsize\bfseries\color{blue!60!black}}{}{0em}{}

\title{
    \textbf{Meridian} \\[0.5em]
    \large Mathematische Grundlagen und Formelreferenz \\[0.3em]
    \normalsize Monte-Carlo-basierte Portfoliorisikoanalyse
}
\author{}
\date{\today}

\begin{document}

\maketitle
\tableofcontents
\newpage

% ─────────────────────────────────────────────────────────────────────────────
\section{Überblick}

Meridian simuliert die zukünftige Wertentwicklung eines Portfolios bestehend aus
mehreren Assets (Aktien, Kryptowährungen) mittels einer \textbf{Monte-Carlo-Simulation}
auf Basis der geometrischen Brownschen Bewegung (GBM). Für jedes Asset werden
historische Marktdaten verwendet, um die statistischen Parameter ($\mu$, $\sigma$,
Korrelationsmatrix) zu schätzen. Die simulierten Endwerte des Portfolios werden
genutzt, um Risikomaße wie VaR und CVaR zu berechnen.

% ─────────────────────────────────────────────────────────────────────────────
\section{Renditeberechnung}

\subsection{Log-Rendite}

Die tägliche Log-Rendite eines Assets wird definiert als:

\begin{equation}
    r_t = \ln\!\left(\frac{P_t}{P_{t-1}}\right)
\end{equation}

\textbf{Erklärung:} Log-Renditen sind symmetrisch um Null und additiv über
mehrere Perioden (im Gegensatz zu einfachen prozentualen Renditen). Sie
entsprechen der stetigen Verzinsung und sind die natürliche Wahl für GBM-basierte
Modelle, da GBM von normalverteilten Log-Renditen ausgeht.

\subsection{Mittlere Log-Rendite}

\begin{equation}
    \mu = \frac{1}{n-1} \sum_{t=1}^{n-1} r_t
\end{equation}

\textbf{Erklärung:} Der Schätzer für den Erwartungswert der täglichen Log-Rendite.
Er wird als täglicher Drift-Parameter $\mu$ in das GBM eingesetzt.

% ─────────────────────────────────────────────────────────────────────────────
\section{Volatilität}

\subsection{Historische Volatilität}

\begin{equation}
    \sigma = \sqrt{\frac{1}{n-1} \sum_{t=1}^{n-1} \left(r_t - \mu\right)^2}
\end{equation}

\textbf{Erklärung:} Die Volatilität ist die Standardabweichung der täglichen
Log-Renditen und ein Maß für das Risiko eines Assets. Ein größeres $\sigma$
bedeutet stärkere Schwankungen des Preises. Der Nenner $n-1$ anstatt $n$
liefert den unverzerrten Schätzer (Bessel-Korrektur).

% ─────────────────────────────────────────────────────────────────────────────
\section{Kovarianzmatrix und Korrelation}

\subsection{Kovarianz}

Die Kovarianz zwischen zwei Assets $i$ und $j$ misst, wie stark ihre Renditen
gemeinsam schwanken:

\begin{equation}
    \text{Cov}(r^{(i)},\, r^{(j)}) = \frac{1}{n-1} \sum_{k=1}^{n}
    \left(r^{(i)}_k - \mu_i\right)\!\left(r^{(j)}_k - \mu_j\right)
\end{equation}

Wichtig: Die Summe läuft nur über \textbf{Tage, an denen beide Assets gehandelt
wurden} (zeitlich ausgerichtete Paare). Mittelwerte $\mu_i$, $\mu_j$ werden
ebenfalls aus dieser Teilmenge berechnet.

\textbf{Erklärung:}
\begin{itemize}[leftmargin=1.5em]
    \item $\text{Cov} > 0$: Assets bewegen sich tendenziell in die gleiche Richtung.
    \item $\text{Cov} < 0$: Assets bewegen sich tendenziell entgegen.
    \item $\text{Cov} = 0$: Keine lineare Abhängigkeit.
\end{itemize}

\subsection{Kovarianzmatrix}

Für ein Portfolio aus $m$ Assets ergibt sich die \textbf{symmetrische, positiv
semidefinite} Kovarianzmatrix:

\begin{equation}
    \boldsymbol{\Sigma} = \begin{pmatrix}
        \text{Cov}(r^{(1)},r^{(1)}) & \cdots & \text{Cov}(r^{(1)},r^{(m)}) \\
        \vdots & \ddots & \vdots \\
        \text{Cov}(r^{(m)},r^{(1)}) & \cdots & \text{Cov}(r^{(m)},r^{(m)})
    \end{pmatrix}
\end{equation}

Die Diagonalelement $\Sigma_{ii} = \text{Var}(r^{(i)}) = \sigma_i^2$ ist die
Varianz des $i$-ten Assets.

\subsection{Normalisierung zur Korrelationsmatrix}

Die Kovarianzmatrix wird auf eine \textbf{Korrelationsmatrix} $\mathbf{P}$
normalisiert:

\begin{equation}
    \rho_{ij} = \frac{\Sigma_{ij}}{\sqrt{\Sigma_{ii} \cdot \Sigma_{jj}}}
    = \frac{\text{Cov}(r^{(i)},r^{(j)})}{\sigma_i \cdot \sigma_j}
\end{equation}

\textbf{Erklärung:} Der Korrelationskoeffizient $\rho_{ij} \in [-1, 1]$ ist
dimensionslos und erlaubt den direkten Vergleich zwischen Asset-Paaren unabhängig
von der absoluten Preishöhe. $\rho_{ij} = 1$: perfekte positive Korrelation;
$\rho_{ij} = -1$: perfekte negative Korrelation.

\subsection{Cholesky-Zerlegung}

Die normalisierte Kovarianzmatrix (Korrelationsmatrix) wird Cholesky-zerlegt:

\begin{equation}
    \mathbf{P} = \mathbf{L} \mathbf{L}^\top
\end{equation}

wobei $\mathbf{L}$ eine untere Dreiecksmatrix ist.

\textbf{Erklärung:} Die Cholesky-Zerlegung wird verwendet, um aus unabhängigen
Zufallsvariablen \textbf{korrelierte} Zufallsvariablen zu erzeugen. Wenn
$\mathbf{Z}$ ein Vektor unabhängiger Zufallsvariablen mit Varianz 1 ist, dann
hat $\boldsymbol{\varepsilon} = \mathbf{L} \mathbf{Z}$ die gewünschte
Kovarianzstruktur $\mathbf{P}$. Die Zerlegung existiert genau dann, wenn
$\mathbf{P}$ positiv definit ist.

% ─────────────────────────────────────────────────────────────────────────────
\section{Schwanzverteilung: Excess Kurtosis und Student-t}

\subsection{Excess Kurtosis}

Die \textbf{Excess Kurtosis} (Exzess-Wölbung) misst, wie „fett" die
Verteilungsenden (Tails) im Vergleich zu einer Normalverteilung sind:

\begin{equation}
    \kappa = \frac{1}{n-1} \sum_{t=1}^{n-1} \left(\frac{r_t - \mu}{\sigma}\right)^4 - 3
\end{equation}

\textbf{Erklärung:}
\begin{itemize}[leftmargin=1.5em]
    \item $\kappa = 0$: Normalverteilung (Mesokurtisch).
    \item $\kappa > 0$: Leptokurtisch. Fettere Tails — extreme Ereignisse treten
          häufiger auf als in der Normalverteilung. Typisch für Finanzzeitreihen.
    \item $\kappa < 0$: Platykurtisch. Dünnere Tails.
\end{itemize}

\subsection{Freiheitsgrade der Student-t-Verteilung}

Aus der Excess Kurtosis werden die Freiheitsgrade $\nu$ der Student-t-Verteilung
geschätzt:

\begin{equation}
    \nu = 4 + \frac{6}{\kappa}, \qquad \kappa > 0
\end{equation}

Ist $4 + 6/\kappa \leq 0$ (d.\,h.\ $\kappa$ sehr klein oder negativ), wird
$\nu \to \infty$ gesetzt, was einer Normalverteilung entspricht.

\textbf{Erklärung:} Die Student-t-Verteilung mit $\nu$ Freiheitsgraden hat fettere
Tails als die Normalverteilung. Sie modelliert besser die empirisch beobachtbare
Häufung extremer Renditen (Crashs, Sprünge). Für $\nu \to \infty$ konvergiert
sie zur Normalverteilung.

% ─────────────────────────────────────────────────────────────────────────────
\section{Monte-Carlo-Simulation}

\subsection{Geometrische Brownsche Bewegung (GBM)}

Jeder Asset-Preis $S_j$ wird täglich durch den folgenden diskreten GBM-Schritt
fortgeschrieben:

\begin{equation}
    S_j(t + \Delta t) = S_j(t) \cdot \exp\!\left[
        \left(\mu_j - \frac{\sigma_j^2}{2}\right)\Delta t
        \;+\;
        \sigma_j \cdot \varepsilon_j \cdot \sqrt{\Delta t}
    \right]
\end{equation}

mit $\Delta t = \frac{1}{252}$ (ein Handelstag in Jahren, da 252 Handelstage
pro Jahr).

\textbf{Erklärung:}
\begin{itemize}[leftmargin=1.5em]
    \item $\mu_j - \frac{\sigma_j^2}{2}$: \textbf{Itô-Korrektur} (auch Drift-Korrektur).
          Die GBM ist im Log-Raum definiert; die $-\sigma^2/2$-Korrektur stellt
          sicher, dass der Erwartungswert von $S$ korrekt ist (Itô-Lemma).
    \item $\sigma_j \cdot \varepsilon_j \cdot \sqrt{\Delta t}$: Stochastischer Term.
          Die Streuung des Preises wächst proportional zu $\sqrt{\Delta t}$
          (Diffusionseigenschaft der Brownschen Bewegung).
\end{itemize}

\subsection{Korrelierte Zufallsvariablen}

Die asset-spezifischen Schocks $\varepsilon_j$ werden korreliert erzeugt:

\begin{equation}
    \boldsymbol{\varepsilon} = \mathbf{L} \cdot \mathbf{Z},
    \qquad Z_k \sim t_{\nu_k}(0, 1)
\end{equation}

wobei:
\begin{itemize}[leftmargin=1.5em]
    \item $\mathbf{Z}$: Vektor unabhängiger, standardisierter Student-t-Zufallsvariablen,
          je Asset mit individuellen Freiheitsgraden $\nu_k$.
    \item $\mathbf{L}$: Untere Cholesky-Dreiecksmatrix der Korrelationsmatrix.
    \item $\boldsymbol{\varepsilon}$: Resultierender Vektor mit der gewünschten
          Korrelationsstruktur $\text{Cov}(\boldsymbol{\varepsilon}) = \mathbf{L}\mathbf{L}^\top = \mathbf{P}$.
\end{itemize}

\subsection{Portfoliowert}

Der Portfoliowert am Ende eines Simulationspfades ist die gewichtete Summe der
Einzelkurse:

\begin{equation}
    V = \sum_{j=1}^{m} w_j \cdot S_j
    = \mathbf{w}^\top \mathbf{S}
\end{equation}

% ─────────────────────────────────────────────────────────────────────────────
\section{Risikomaße}

\subsection{Value at Risk (VaR)}

Der \textbf{Value at Risk} zum Konfidenzniveau $(1-\alpha)$ ist das
$\alpha$-Quantil der Verteilung der Portfolioendwerte:

\begin{equation}
    \text{VaR}_\alpha = F_V^{-1}(\alpha)
\end{equation}

In der Simulation wird das $\alpha$-Quantil empirisch über die sortierten
Portfolioendwerte bestimmt:

\begin{equation}
    \text{VaR}_\alpha \approx \text{sorted}(V)\!\left[\lfloor \alpha \cdot N \rfloor\right]
\end{equation}

\textbf{Erklärung:} $\text{VaR}_{0{,}05}$ bedeutet: Mit 95\,\% Wahrscheinlichkeit
liegt der Portfoliowert nach dem Simulationshorizont \emph{über} diesem Wert.
Anders gesagt: Im schlechtesten 5\,\% der Fälle fällt der Wert unter $\text{VaR}_{0{,}05}$.

\subsection{Conditional Value at Risk (CVaR) / Expected Shortfall}

Der \textbf{CVaR} (auch Expected Shortfall, ES) ist der \emph{Erwartungswert}
des Portfoliowertes, \emph{gegeben} dass er unter dem VaR liegt:

\begin{equation}
    \text{CVaR}_\alpha = \mathbb{E}\!\left[V \,\Big|\, V \leq \text{VaR}_\alpha\right]
    \approx \frac{1}{\lfloor \alpha N \rfloor + 1}
    \sum_{k=0}^{\lfloor \alpha N \rfloor} \text{sorted}(V)[k]
\end{equation}

\textbf{Erklärung:} CVaR ist ein kohärenteres Risikomaß als VaR, da es nicht
nur fragt „ab wann wird es schlimm?" sondern auch „wie schlimm wird es im
Durchschnitt, wenn es schlimm wird?". Er berücksichtigt die Form der Verteilung
im Tail.

% ─────────────────────────────────────────────────────────────────────────────
\section{Glossar der Variablen}

\renewcommand{\arraystretch}{1.4}
\begin{tabular}{>{\bfseries}p{2.5cm} p{3.5cm} p{8cm}}
    \toprule
    \textbf{Symbol} & \textbf{Einheit / Typ} & \textbf{Bedeutung} \\
    \midrule
    $P_t$       & Preis [USD]         & Bereinigter Schlusskurs (\emph{adjusted close}) des Assets zum Zeitpunkt $t$ \\
    $r_t$       & dimensionslos       & Log-Rendite von $t-1$ nach $t$: $\ln(P_t / P_{t-1})$ \\
    $n$         & ganze Zahl          & Anzahl der historischen Datenpunkte \\
    $\mu$       & 1/Tag               & Mittlere tägliche Log-Rendite (Drift-Parameter) \\
    $\sigma$    & 1/$\sqrt{\text{Tag}}$ & Tägliche Volatilität (Standardabweichung der Log-Renditen) \\
    $\kappa$    & dimensionslos       & Excess Kurtosis der Log-Renditenverteilung \\
    $\nu$       & dimensionslos       & Freiheitsgrade der Student-t-Verteilung \\
    $\boldsymbol{\Sigma}$ & Matrix    & Kovarianzmatrix der Log-Renditen aller Assets \\
    $\mathbf{P}$ & Matrix             & Normalisierte Kovarianzmatrix (Korrelationsmatrix), $\rho_{ij} \in [-1,1]$ \\
    $\rho_{ij}$ & dimensionslos       & Korrelationskoeffizient zwischen Asset $i$ und Asset $j$ \\
    $\mathbf{L}$ & untere Dreiecksmat.& Cholesky-Faktor: $\mathbf{P} = \mathbf{L}\mathbf{L}^\top$ \\
    $\mathbf{Z}$ & Vektor             & Vektor unabhängiger Student-t-Zufallsziehungen \\
    $\boldsymbol{\varepsilon}$ & Vektor & Korrelierter Zufallsvektor: $\mathbf{L}\mathbf{Z}$ \\
    $S_j(t)$    & Preis [USD]         & Simulierter Kurs von Asset $j$ zum Zeitpunkt $t$ \\
    $\Delta t$  & Jahre               & Zeitschritt der Simulation: $1/252$ \\
    $m$         & ganze Zahl          & Anzahl der Assets im Portfolio \\
    $w_j$       & dimensionslos       & Gewicht von Asset $j$ im Portfolio ($\sum w_j = 1$) \\
    $V$         & Preis [USD]         & Gesamter Portfoliowert am Ende des Simulationshorizonts \\
    $N$         & ganze Zahl          & Anzahl der Monte-Carlo-Simulationspfade \\
    $\alpha$    & dimensionslos       & Signifikanzniveau für VaR/CVaR (z.\,B.\ $0{,}05$ für 95\,\%-VaR) \\
    $\text{VaR}_\alpha$ & Preis [USD] & Value at Risk: $\alpha$-Quantil der simulierten Portfolioendwerte \\
    $\text{CVaR}_\alpha$ & Preis [USD] & Conditional VaR: Erwartungswert der Portfolioendwerte unterhalb $\text{VaR}_\alpha$ \\
    \bottomrule
\end{tabular}

\end{document}
